\documentclass[14pt]{extarticle}
\usepackage{fontspec}
\usepackage{polyglossia}
\usepackage{amsmath}
\usepackage{amsfonts}
\usepackage{amssymb}
\usepackage{pgfplots}
\pgfplotsset{compat=1.18}

% Hebrew setup
\setdefaultlanguage{hebrew}
\setotherlanguage{english}
\newfontfamily\hebrewfont[Script=Hebrew]{Arial}

% Math font setup for proper Greek symbols
\usepackage{unicode-math}
\setmathfont{Latin Modern Math}

% Increase line spacing throughout document
\linespread{1.3}

\usepackage[margin=1in]{geometry}

% Left-aligned equation environment
\newenvironment{lefteqs}{%
    \begin{english}%
    \begin{flushleft}%
    \renewcommand{\arraystretch}{1.3}%
    $\begin{array}{@{}l@{}}%
}{%
    \end{array}$%
    \end{flushleft}%
    \end{english}%
}

\begin{document}

\begin{center}
    {\LARGE \textbf{מטלה 1 - אלקטרומגנטיות אנליטית}}\\[0.5em]
    {\large שי רואימי - 318986270}\\[0.3em]
\end{center}

\hrule
\vspace{2em}

\textbf{\large שאלה 1}

(שאלה 4 בפרק המבוא)
נתון $\phi=x^2yz+4xz^2$

\textbf{א)}
נחשב את $\vec{\nabla} \phi$:

\begin{lefteqs}
    \vec{\nabla} \phi = \frac{\partial \phi}{\partial x} \hat{x} + \frac{\partial \phi}{\partial y} \hat{y} + \frac{\partial \phi}{\partial z} \hat{z}
    = (2xyz+4z^2)\hat{x} + x^2z\hat{y} + (x^2y+8xz)\hat{z}
\end{lefteqs}

\textbf{ב)}
נחשב את הנגזרת הכיוונית של $\phi$ בנקודה $(1,-2,-1)$ בכיוון $2\hat{x}-\hat{y}-2\hat{z}$

וקטור הכיוון הוא $(2,-1,-2)$ וגודלו $\sqrt{2^2+(-1)^2+(-2)^2}=\sqrt{9}=3$,
ולכן וקטור היחידה בכיוון זה הוא $\hat{u}=\frac{1}{3}(2,-1,-2)$.

ערכו של $\vec{\nabla} \phi$ בנקודה $(1,-2,-1)$ הוא: $8\hat{x} - \hat{y} -10\hat{z}$.

\begin{lefteqs}
    \vec{\nabla} \phi \cdot \hat{u} = \frac{1}{3} (8,-1,-10) \cdot (2,-1,-2) = \frac{1}{3} (16 + 1 +20) = \frac{37}{3}
\end{lefteqs}
ערכה של הנגזרת הכיוונית הוא $\frac{37}{3}$.

\vspace{3em}
\hrule
\vspace{2em}

\textbf{\large שאלה 2}

(שאלה 12 בפרק המבוא)
נתון צילינדר שרדיוסו 2 וגובהו 5 כאשר ציר הסימטריה שלו מתלכד עם ציר ה-$z$, והדופן התחתונה שלו במישור $xy$.

בהנתן $\vec{F}=\vec{\nabla} (x^2+y^2)$

\textbf{א)}
נמצא את $\int\int\int \vec{\nabla} \vec{F} \, dV$

\begin{lefteqs}
    \vec{F} = 2x \hat{x} + 2y \hat{y} \\
    \vec{\nabla} \cdot \vec{F} = \frac{\partial (2x)}{\partial x} + \frac{\partial (2y)}{\partial y} = 2 + 2 = 4 \\
    \int\int\int 4 \, dV = 4 \cdot \pi R^2 h = 4 \cdot \pi \cdot 4 \cdot 5 = 80 \pi
\end{lefteqs}

\textbf{ב)}
נמצא את השטף הכולל דרך דפנות הצילינדר

על פי משפט הדיביגרנץ, השטף דרך כל הגליל שווה לתוצאה בסעיף א', כלומר $\phi = \int\int\int (\vec{\nabla} \cdot \vec{F}) \, dV = 80 \pi$.

השטף דרך הפאה העליונה והתחתונה שווה לאפס כיוון של-$\vec{F}$ אין רכיב בכיוון $z$, 
והנומרל בשתי הפאות הוא $\hat{z}$ ו $-\hat{z}$ בהתאמה, ולכן $\vec{F} \cdot d\vec{S} = 0$.

לכן השטף דרך הדפנות הוא $80 \pi$.

\textbf{ג)}
נאשר את קיומו של משפט הדיבגרנץ: כלומר נראה ש $\int\int \vec{F} \cdot d\vec{S} = 80 \pi$

ראשית נזכיר כפי שהוסבר בתשובה לסעיף ב' שהשטף דרך הפאות הוא אפס, ולכן נחשב רק את השטף דרך הדפנות. וכן:

\begin{lefteqs}
    \vec{F} = 2x \hat{x} + 2y \hat{y} = 2(x \hat{x} + y \hat{y}) = 2r \hat{r} \\
    \phi = \int \int \vec{F} \cdot d\vec{S} = 4 \int \int dS = 4 \cdot 2\pi r \cdot h = 80 \pi
\end{lefteqs}

קיבלנו תשובה זהה לסעיף א', ולכן אישרנו את קיומו של משפט הדיבגרנץ.


% \begin{english}
%     \begin{flalign*}
%         (1) \quad & \Psi(x,t)=A \cdot e^{-(kx-\omega t)^2}       & \\
%         (2) \quad & \Psi(x,t)=A \cdot e^{-a^2x^2+i(kx-\omega t)} &
%     \end{flalign*}
% \end{english}

% כאשר
% $J(x,t)=\frac{i \hbar}{2m}(\Psi \frac{\partial \Psi^*}{\partial x} - \Psi^* \frac{\partial \Psi}{\partial x})$
% ,ו $\rho(x,t)=|\Psi(x,t)|^2$.

% עבור (1):

% \begin{lefteqs}
%     \rho(x,t) = A^2 e^{-2(kx-\omega t)^2} \\
%     \dfrac{\partial \Psi}{\partial x} = \dfrac{\partial \Psi^*}{\partial x} \\
%     \Psi=\Psi^* \\
%     J(x,t) = 0
% \end{lefteqs}


\end{document}